\scanpage{58}
celle associée à l'autom suivant de $(I\xrightarrow{p}\Lambda)$ : son action sur $\Lambda$ est la permutation qui fixe $\alpha$ i.e. échange $\beta,\gamma$, son action sur $I$ est l'identité sur $I_\alpha$, et elle échange en\unsure{tre} eux $x_\beta\in I_\beta$ et $\varphi(x_\beta)$. On voit immédiatement que le groupe engendré par ces opérations est le groupe de tous les autom de $\Delta$ laissant invariant sa partition $\Delta=\bigcup X_\alpha$.

\deleted{Le groupe engendré} Le sous-groupe engendré par les $\sigma_{\underline t}$ avec $t$ \unsure{triang.|triangle} de $X_\alpha$ ($\alpha$ fixé) est formé des autom de $\Delta$ qui a) \unsure{invarient} $X_\alpha$ (donc la partition)
b) opèrent trivialement sur l'ens $\unsure{\tau|=|\pi}(X_\alpha)\simeq I_\alpha$ des ``trois orbites'' associé à $X_\alpha$.

\textcircled{16} Soit $X\subset S$ un bitriangle de $S$, et $\deleted{\unsure{/}}(S_i)_{i\in I_X}$ l'une des ``trois orbites'' de $S-X$.

Alors les sous-graphes \inserted{de type} $\unsure{\mathscr T|\Gamma|\mathcal T}_1$ \deleted{\unsure{$\subset$}} contenant $X$ sont les $X\cup S_i$. Donc la \deleted{\unsure{donnée}} \inserted{\unsure{catég\textsuperscript{rie}|catégorie}} \deleted{\unsure{de leurs}} \unsure{des} triplets $(S\supset\unsure{\mathscr X|\mathfrak X|\mathcal X}\supset X)$ d'un \deleted{\unsure{cubique}} \inserted{graphe}

% archive sheet number: 57

